
% Introdução até 18/05/2026
From an intertemporal social-efficiency perspective, neither extreme predation, defined by the indiscriminate and often illegal depletion of natural resources, nor strict preservationism, characterized by minimal or non-extractive use of the environment, can provide a stable and efficient resolution of the trade-off between production and ecological integrity \cite{dasgupta2001human}. These two poles represent opposing boundaries of environmental interaction: the former disregards sustainability entirely, while the latter may overly constrain productive potential under conventional productivity metrics, which typically emphasize output expansion over ecological feedbacks. Between these extremes lies the economic problem of resource extraction: the allocation of renewable and non-renewable natural resources over time, under conditions of scarcity, regeneration, depletion, opportunity costs, and intertemporal trade-offs \cite{hotelling1931economics, clark2010mathematical, dasgupta2001human}. In this narrower economic sense, extraction need not be predatory; it may remain compatible with sustainability when the rate of use respects regenerative capacity, stock constraints, and the shadow value of future resource availability.

The concept of extractivism, however, carries a different meaning in the political economy literature. Rather than denoting resource extraction as an intertemporal allocation problem, extractivism refers to a development pattern based on the large-scale appropriation of natural resources, frequently oriented toward commodity exports, limited domestic value addition, and ecological degradation \cite{gudynas2009diez, acosta2013extractivism, mckay2017agrarian}. Under this interpretation, the issue is not extraction per se, but the institutional and productive regime through which extraction becomes structurally tied to depletion, dependency, and asymmetric distribution of environmental costs.

% ---

